% freewill_note.tex — companion note to unification.tex
% RUN provenance:
%   run 1 (2026-07-18): drafted from a four-reader extraction sweep of unification.tex.
%   run 2 (2026-07-18): adversarially certified by a 110-agent swarm (6 attack lenses ->
%     per-finding refuter votes -> completeness critic). 31 confirmed findings (3 FATAL)
%     and 2 coverage gaps folded into this version; 31 candidate findings refuted.
%     Key repairs: F2a scoped (protocol admissibility constraints exist: ax:complete,
%     (RP)); F3b restated time-indexed (K IS record-measurable under ax:complete);
%     F1b corner re-tensed (jailer = terminal record, a future-boundary datum, and the
%     BRANCH's records, not the being's own); F1b/F1c regraded onto the net bundle;
%     regime scoping added (no single regime carries the whole conjunction; conj:tomita
%     named); corollary (ii) scoped to the passive law; inter-being influence conceded.
%   run 3 (2026-07-18): budgeted 8-agent targeted verification of the run-2 fix
%     sites (one pinned checker per cluster; 250k tokens). All 8 clusters verified
%     faithful; 2 MINOR wording slips found and fixed in this version (F2a final
%     clause overclaimed record-generation; rem:regimes silently upgraded F2c's
%     regime). No FATAL/MAJOR residue.
% Grade discipline: grades carried faithfully; no PLAUSIBLE consumed at THEOREM;
%   every clause carries its own conditionality tag and its regime.
% ZERO writes to unification.tex / ISSUES.md. Self-contained; compiles standalone.
% Parent-paper cross-references are given by NAME in prose with the exact \label
%   + line anchor in a %-comment, so this file has ZERO undefined refs.
\documentclass[11pt]{article}
\usepackage[margin=1in]{geometry}
\usepackage[T1]{fontenc}
\usepackage{lmodern}
\usepackage{amsmath,amssymb,amsthm}
\usepackage{booktabs}
\usepackage{array}
\usepackage{xcolor}
\usepackage[colorlinks=true,linkcolor=blue!50!black,citecolor=blue!50!black]{hyperref}

\emergencystretch=1.5em

\newtheorem{definition}{Definition}
\newtheorem{lemma}{Lemma}
\newtheorem{theorem}{Theorem}
\newtheorem{corollary}{Corollary}
\newtheorem{proposition}{Proposition}
\newtheorem{remark}{Remark}

% status macros (tbrme-uniqueness-outline house style)
\newcommand{\THM}{{\normalfont\scshape theorem}}                % theorem-grade in its stated regime
\newcommand{\THMC}[1]{{\normalfont\scshape theorem}$\,\mid\,${\normalfont\small #1}} % conditional; bundle listed
\newcommand{\ADOPTED}{{\normalfont\scshape adopted}}            % ambient-literature import
\newcommand{\STRUCT}{{\normalfont\scshape structural}}          % true by inspection of the axiom set
\newcommand{\DERIVED}{{\normalfont\scshape derived}}            % one-step corollary proved in this note
\newcommand{\CONCEDED}{{\normalfont\scshape conceded}}          % explicitly not claimed

\newcommand{\MM}{\mathcal M}
\newcommand{\HH}{\mathcal H}
\newcommand{\AAA}{\mathcal A}
\newcommand{\KK}{\mathcal K}
\newcommand{\FF}{\mathcal F}
\newcommand{\PP}{\mathcal P}
\newcommand{\E}{\mathbb E}
\newcommand{\Prob}{\mathbb P}
\newcommand{\Tr}{\operatorname{Tr}}
\newcommand{\id}{\mathbf 1}
\newcommand{\RC}{\mathrm{RC}}

\title{Operational freedom in the two-boundary retrocausal framework:\\
an equal-possibility theorem\\
\large (companion note --- hospitality and equality, never metaphysics)}
\author{Alexandru Dima}
\date{}

\begin{document}
\maketitle

\begin{center}
\fbox{\parbox{0.92\textwidth}{\small
\textbf{Metaphysics fence (binding).} This note does \emph{not} and \emph{cannot} prove
libertarian (metaphysical) free will. What it proves is that a precisely defined
\emph{operational} surrogate --- the conjunction of four formal properties
(equal possibility, non-coercion, openness, efficacy) --- holds for every admissible
observer of the parent framework, with the same strength for every observer, at the
regime- and bundle-scoped grades stated clause by clause.
Whether that conjunction deserves the name ``free will'' is a philosophical
judgment outside physics; the note's own position is compatibilist and is argued,
not smuggled. The parent paper defines an observer operationally (ideal-QRF axioms)
and explicitly disclaims any appeal to consciousness or agency; this note inherits
that disclaimer verbatim. % parent: prop:observer-invariance, tier (C) (unification.tex l.23957)
Deliberation, intention, and phenomenology are not modeled anywhere in the framework
and are not modeled here.}}
\end{center}

\begin{abstract}
\noindent Within the two-boundary retrocausal framework of the parent paper, an
observer (``being'') is a worldline--protocol--filtration triple. We define, in the
framework's own vocabulary, four properties that jointly constitute
\emph{operational freedom}: (F1) \emph{equal possibility} --- at every moment the
algebra of available instruments is the same, up to $*$-isomorphism, for every
being; no available option is assigned probability zero by the two-boundary
conditioning except through incompatibility with the branch's terminal record;
and every being has the same (dense) operational reach; (F2) \emph{non-coercion}
--- no axiom determines or probabilistically weights the choice of protocol, no
spacelike intervention by any other being can alter the being's marginal
statistics, and the final boundary condition is operationally inert in prior
average; (F3) \emph{openness} --- the being's forecast of the terminal record is
an unbiased martingale, no variable available before the terminal limit predicts
it better, and definiteness arrives only with the records themselves; (F4)
\emph{efficacy} --- the being's choices shift the law of the realized branch,
lawfully and causally. The main theorem asserts F1--F4 for every admissible
being, clause by clause, each at its honest grade and regime: the Part-I spine
clauses are theorems of the finite-dimensional core; the field-theoretic clauses
carry a named net bundle; the measurement-independence and no-better-predictor
clauses are conditional on the parent's \textsf{CHT-POVM} conjecture plus the
thermal-time postulate; and the transfer of the whole conjunction to a single
field-theoretic being is scoped by the parent's type-III lift (Araki branch
theorem-grade; sandwich branch \textsf{conj:tomita}). The corollary that gives
the note its title: \emph{at any moment, all beings have the same possibilities
of choice} --- isomorphic choice algebras, a common passive-law terminal
distribution, universal reach, common outer time flow (postulate-grade) ---
while what differs between beings (location, records, thermodynamic state) is
exhaustively indexical. An attack ledger and a grade table close the note.
\end{abstract}

\section{Setting and vocabulary}\label{sec:setting}

All objects are the parent paper's. $(\MM,g)$ is a fixed globally hyperbolic
background with Cauchy temporal function $\tau$ and slab $\MM_{[\tau_i,\tau_f]}$;
$\rho_{\rm i}$ is the initial state on $\Sigma_{\tau_i}$ and $\{E_k\}_{k\in\KK}$
the terminal POVM on $\Sigma_{\tau_f}$, with backward effects
$\Lambda_k(\tau)$ and time-invariant Born weights $p(k)=\Tr(\rho(\tau)\Lambda_k(\tau))$.
% parent: ax:data (unification.tex l.583); def:rho2 (l.817); lem:born (l.754)
The conditioned state is the symmetric-sandwich branch
$\rho_2(\tau\mid k)=\rho(\tau)^{1/2}\Lambda_k(\tau)\rho(\tau)^{1/2}/p(k)$, unitarily
covariant with the same generator $H(\tau)$.
% parent: thm:unitary (l.942)
Local algebras $\AAA(R)$ satisfy isotony, microcausality (M2) and causal
propagation (M3).
% parent: ax:microcausality (l.589)

Two probability laws must be kept apart throughout, as the parent insists: the
\emph{passive} law $\Prob_\gamma$, defined by construction so that its terminal
marginal is $p(k)$ for every protocol, and the \emph{active} sequential-Born law
$P_M$, describing physically performed interventions, whose terminal marginal is
the disturbed weight $\widetilde p(k)$ and need not equal $p(k)$.
% parent: passive law and by-construction marginal (l.1188--1191); active law
% distinction (l.1229--1241); disturbed Born law fd:thm:instrument (iv) (l.14678)
Statements of ``sameness of law'' below are passive-law statements; statements of
``efficacy'' are active-law statements. Conflating the two was the single most
frequent defect found by this note's adversarial review, and the scoping is now
explicit at every use.

\emph{Regimes.} The parent proves its Part-I spine in the finite-dimensional
core (all operator identities literal), and its conditioning, Reeh--Schlieder,
and no-go--evasion results in the field-theoretic regime (type III$_1$ net).
The type-III lift of the sandwich-state machinery is the parent's conjecture
\textsf{conj:tomita}; its Araki-perturbation branch is realised at theorem grade
(the parent's (E4) is closed unconditionally), while the collapse chain
transfers type-free through the algebra-free predual theorem.
% parent: conj:tomita (l.4699); Araki realization / E4 closure per app. G and the
% sharpened summary; fd:thm:predual (l.14898)
Each clause below names the regime in which its grade is earned; the transfer of
the full conjunction to one field-theoretic being is scoped in
Remark~\ref{rem:regimes}.

\begin{definition}[Being; admissibility]\label{def:being}
A \emph{being} is the parent paper's observer datum: a triple
$B=(\gamma,\PP_\gamma,\FF^\gamma_\bullet)$ consisting of an inextendible
future-directed timelike worldline $\gamma$ crossing each slice once, a fixed
measurement protocol $\PP_\gamma=(\{M^{(m)}_j\}_j)_{m}$ with
$M^{(m)}_j\in\AAA(D_{\varepsilon_m}(\gamma(s_m)))$ for some resolutions
$\varepsilon_m>0$, and the record filtration $\overline\FF^\gamma_\tau$ it
generates under the passive law $\Prob_\gamma$.
% parent: sec:filt, worldline+protocol+filtration definitions (l.1128--1246);
% the parent's point-algebra formulation A(gamma(s_m)) is the outer-regularized
% finite-dimensional core version; the diamond formulation used here specializes
% to it there and is the non-trivial one in the field-theoretic regime, where
% point algebras are trivial. The parent's Part-I formulas use only that the
% M^{(m)}_j lie in a local algebra at the slice, so they go through verbatim.
A being is \emph{admissible} for a result if it satisfies the protocol
conditions that result assumes: informativeness where terminal
record-completeness is invoked, and the record-layer restriction (RP) where the
record-sourced geometry is invoked. These are genuine admissibility conditions
of the parent, not conveniences of this note.
% parent: ax:complete (l.1654); (RP) ``a genuine restriction on P_gamma, not
% bookkeeping'' (l.2445--2449); uniformly informative protocol class (l.3603--3605)
A \emph{branch} $k$ is admissible if $p(k)>0$. When frame-invariance of
predictions is invoked, $B$ is additionally assumed to satisfy the ideal-QRF
axioms of the parent observer appendix.
% parent: prop:observer-invariance, tier (C) (l.23957)
\end{definition}

\begin{definition}[Choice situation]\label{def:choice}
For a being $B$, a moment $\tau$, and a resolution $\varepsilon>0$, the
\emph{choice set} $\mathsf C_B(\tau;\varepsilon)$ is the set of all finite Kraus
instruments $\{M_j\}$ with $M_j\in\AAA(D_\varepsilon(\gamma(\tau)))$ and
$\sum_j M_j^{*}M_j=\id$, where $D_\varepsilon(x)$ is the causal diamond of radius
$\varepsilon$ about $x$. By construction, every step of an admissible protocol
$\PP_\gamma$ is an element of the choice set at its moment and resolution.
\end{definition}

\begin{lemma}[Functoriality of choice sets]\label{lem:functorial}
Any unital $*$-isomorphism of local algebras carries choice sets bijectively
onto choice sets.
\end{lemma}
\begin{proof}
A unital $*$-isomorphism preserves adjoints, sums, and the unit, hence the
Kraus normalization $\sum_j M_j^*M_j=\id$; its inverse does the same.
\end{proof}

\begin{definition}[Operational freedom]\label{def:freedom}
A being $B$ is \emph{operationally free} if the following four properties hold.
\begin{itemize}
\item[\textbf{(F1)}] \emph{Equal possibility.}
  (a) For any two beings $B,B'$, moments $\tau,\tau'$, and resolutions
  $\varepsilon,\varepsilon'$ in the field-theoretic regime, the choice sets
  $\mathsf C_B(\tau;\varepsilon)$ and $\mathsf C_{B'}(\tau';\varepsilon')$ are
  identified by a $*$-isomorphism of the underlying local algebras.
  (b) \emph{Non-extinction:} for every admissible branch $k$, no nonzero effect
  available to $B$ is assigned probability zero by the conditioned state, except
  on the corner fixed by the branch's terminal record --- a future-boundary
  datum, block-fixed but in general not yet recorded at $\tau$.
  (c) \emph{Universal reach:} from $B$'s diamond alone, local operations applied
  to the conditioned state reach a dense set of global states.
\item[\textbf{(F2)}] \emph{Non-coercion.}
  (a) No axiom or equation of the framework determines or probabilistically
  weights the protocol $\PP_\gamma$; the framework's protocol conditions
  (informativeness, (RP)) delimit which protocols support which of its
  results, never which protocol the being adopts.
  (b) Measurement independence holds at the initial slice (the framework is not
  superdeterministic).
  (c) No intervention by any spacelike-separated being alters $B$'s marginal
  statistics.
  (d) The final boundary condition is operationally inert for $B$ in prior
  average: averaged over branches, $B$'s statistics are exactly those of forward
  quantum mechanics.
\item[\textbf{(F3)}] \emph{Openness.}
  (a) $B$'s posterior $\pi^\gamma_\tau(k)$ over the terminal record is a
  martingale on $B$'s own filtration, unbiased under every record-measurable
  stopping strategy, and this survives measurement disturbance.
  (b) No variable available strictly before the terminal limit --- and no
  extension of the framework's statistics beyond $\{E_k\}$ --- predicts $K$
  better than the posterior.
  (c) Definiteness arrives only with the records: $\pi^\gamma$ collapses to
  $\mathbf 1_{\{K=k\}}$ only in the terminal limit, and for horizon-restricted
  beings may never fully collapse.
\item[\textbf{(F4)}] \emph{Efficacy.}
  (a) $B$'s choice of instrument changes the law of the realized branch:
  the terminal weights are intervention-dependent under the active law.
  (b) The consequences are lawful: between interventions, conditioned
  expectations obey the conserved Ehrenfest dynamics along the realized branch,
  and the intervention step itself is governed by the instrument update.
  (c) The consequences are causal, not anticipating: a deferred decision enters
  the record-sourced geometry only after its record lies in the causal past of
  the readout.
\end{itemize}
\end{definition}

\section{The theorem}\label{sec:theorem}

\begin{theorem}[Operational freedom of every being]\label{thm:freedom}
Every admissible being of the parent framework is operationally free in the
sense of Definition~\ref{def:freedom}, with the following clause-by-clause
provenance, regime, and conditionality.
\begin{itemize}
\item[\textbf{(F1a)}] \ADOPTED{} $+$ net bundle; field-theoretic regime. The
parent's field-theoretic results assume a Haag--Kastler net of type III$_1$ on
globally hyperbolic spacetime.
% parent: conj:tomita setting (l.4699); thm:reeh-schlieder-ABL hypotheses (l.22739)
Granting in addition that every relatively compact diamond algebra is a
\emph{factor} with separable predual and that the net has the
split/hyperfiniteness property standard for well-behaved nets, every such
algebra is the unique hyperfinite type III$_1$ factor (Haagerup uniqueness;
Connes classification), hence any two beings' diamond algebras --- at any two
moments, at any two resolutions --- are $*$-isomorphic, and by
Lemma~\ref{lem:functorial} their choice sets are identified. \emph{Honest
fences:} (i) factoriality and predual separability are part of the import
bundle, not stated at the parent anchor --- a nontrivial center at any diamond
would break the pairwise identification; (ii) in the finite-dimensional Part~I
core this clause can fail (unital $*$-subalgebras of $\mathcal L(\HH)$ at
different worldline points need not be isomorphic). Equality of possibility is
a genuinely field-theoretic fact, and the type III$_1$ structure is
load-bearing, not decorative.
\item[\textbf{(F1b)}] \THMC{net hypotheses}; field-theoretic regime; dense-range
branch class exact, corner form in general. The parent's Reeh--Schlieder theorem
for ABL conditioning --- proved under its stated hypotheses: type III$_1$
Haag--Kastler net, faithful normal Hadamard state with the Reeh--Schlieder
property (ambient imports), $L_k=\Lambda_k(\tau)^{1/2}$ bounded --- gives: if
$L_k$ has dense range the conditioned state $\omega_k$ is faithful and normal on
every local algebra, so $\omega_k(F)>0$ for every nonzero effect $F$ available
to any being at any moment.
% parent: thm:reeh-schlieder-ABL, clauses (i)--(iii) (l.22739)
If $E_k$ is projective with nontrivial kernel, faithfulness is preserved
exactly on the support corner and lost exactly on $\ker L_k$.
% parent: thm:reeh-schlieder-ABL, clause (iv); rem:reeh-schlieder-status:
% ``the post-selection collapse ... onto outcomes compatible with k'' (l.22834--22838)
The corner is fixed by the \emph{terminal record of the branch}: a
future-boundary datum, block-fixed at every $\tau$ but in general not yet
recorded at $\tau$; under terminal record-completeness it coincides in the
terminal limit with the branch's accumulated record history --- the branch's,
not the individual being's, since other beings' recorded outcomes are part of
it. What the clause rules out is extinction by location, constitution, or by
another being's \emph{protocol choices}; extinction by facts recorded on the
branch (by anyone) is exactly what conditioning is.
\item[\textbf{(F1c)}] \THMC{net hypotheses} $+$ \DERIVED{}; field-theoretic
regime, dense-range class. Cyclicity of the conditioned vector for every local
algebra
% parent: thm:reeh-schlieder-ABL, clause (ii) (l.22739)
plus the standard instrument construction (any $A\in\AAA(D)$ with
$\|A\|\le1$ is a success branch $M_1=A$ of the two-outcome instrument
$\{A,\sqrt{\id-A^*A}\}$; rescale a cyclic-orbit approximant) yields: for every
global normal state $\varphi$ and every $\delta>0$ there is an operation in
$B$'s diamond, of strictly positive success probability, whose normalized
output state is within $\delta$ of $\varphi$ in vector norm on the GNS space.
The derivation from cyclicity is this note's (one step, sketched above); the
cyclicity is the parent's. Every being has the same dense reach; postselected,
and with possibly minute success probability, but available identically to all.
\item[\textbf{(F2a)}] \STRUCT{}; scoped. By inspection of the axiom set: no
axiom or equation determines the protocol or assigns it a probability --- the
parent introduces $\PP_\gamma$ as a fixed non-random external input.
% parent: ``The POVM choices are encoded in P_gamma and are not random'' (l.1150--1151)
What inspection does \emph{not} yield --- and this note's first version wrongly
asserted --- is that no axiom constrains $\PP_\gamma$ at all: terminal
record-completeness is a joint identifiability condition on the pair
$(\PP_\gamma,K)$ that fails for uninformative protocols, and the Part-II
record layer imposes (RP), ``a genuine restriction on $\PP_\gamma$, not
bookkeeping.''
% parent: ax:complete (l.1654); (RP) restriction (l.2445--2449)
These are admissibility conditions delimiting which protocols support
collapse and record-sourcing; they do not weight, forbid, or select the
being's choice among admissible protocols. Hospitality, correctly scoped:
the framework has no resource to force or bias the choice; it does say which
choices yield records that identify the outcome and may source geometry.
\item[\textbf{(F2b)}] \THMC{Conj.\ \textsf{CHT-POVM} $+$ Pos.\ thermal-time}.
The parent's combined Bell-style evasion theorem --- whose standing hypotheses
are the \textsf{CHT-POVM} conjecture, the thermal-time postulate, terminal
record-completeness, and bulk no-signalling --- places the framework as
``non-dynamical (all-at-once), future-input-dependent retrocausal without
operational signalling'' and states verbatim that it is \emph{not}
superdeterministic: measurement independence at the initial slice is retained;
what is abandoned is No-Retrocausality, and the resulting future-input
dependence of past hidden variables is a modular invariant, not a tunable
parameter.
% parent: thm:bell-no-go-evasion, hypotheses l.22993--22995, item (1)
% (meas.-indep. at l.23014--23016); conditionality per rem:bell-no-go-status
% (l.23100), inherited from thm:wood-spekkens-evasion (l.22491) via
% conj:CHT-POVM (l.21659).
The parent's own concessions are inherited unshrunk: the evasion is of
\emph{causal} fine-tuning, not Catani--Leifer \emph{operational} fine-tuning
(``lawful fine-tuning, not its absence''), and a general adaptive-interventions
analysis is an open scope item.
% parent: rem:WS-operational-scope (l.22566)
\item[\textbf{(F2c)}] \THM{}; Part-I core (finite-dimensional). Bulk
no-signalling: for spacelike separated regions, the choice of instrument in one
leaves the outcome \emph{marginals} in the other exactly invariant; the
past-to-future direction uses microcausality and causal propagation, the
future-to-past direction not even that.
% parent: thm:no-signal (l.1340); consolidated statement rem:nosignal-coherent (l.1602)
Scope, stated plainly: this protects $B$'s $k$-summed marginals. It does not
protect per-branch post-selected statistics, and it does not make the shared
terminal law insensitive to others (see Remark~\ref{rem:mutual}).
\item[\textbf{(F2d)}] \THM{} in the Part-I core; in the field-theoretic regime
via the Araki-conditioning branch, with the sandwich-branch lift carried by
\textsf{conj:tomita}. The prior-average identity
$\E_{\Prob_\gamma}[\rho^\gamma_\RC(\tau)]=\rho(\tau)$ and the reduction
corollary: averaged over branches, every expectation any being can form is the
ordinary forward-QM expectation from $\rho_{\rm i}$.
% parent: lem:post-props (iii) (l.1257); cor:reduction (l.1312); type-III
% averaging identity on the Araki branch per thm:E4-dual-weight (iv) as cited in
% the parent's conj:tomita discussion (l.4699ff)
The future boundary shapes which past is realised and provably cannot be used
to send a message into it.
% parent: closing line of rem:nosignal-coherent (l.1602)
\item[\textbf{(F3a)}] \THM{} in the Part-I core; type-free in general
(disturbance-robust) $+$ \DERIVED{}. The posterior is a martingale on the record
filtration by the tower property; the martingale property survives arbitrary
measurement disturbance --- disturbance changes the measure, not the Bayesian
structure --- and the collapse chain is algebra-free (valid verbatim for type
III preduals).
% parent: lem:post-props (i) (l.1257); fd:thm:instrument, clause (ii) (l.14678);
% fd:thm:predual (l.14898)
Derived here (standard probability, one step): for every record-measurable
stopping time $T$, $\E[\pi^\gamma_T(k)]=\pi^\gamma_0(k)$ (optional sampling for
a bounded martingale closed by $\mathbf 1_{\{K=k\}}$). This is
\emph{unbiasedness} under stopping strategies, not non-improvement of accuracy:
expected forecast accuracy does improve as records accumulate --- that is
exactly F3c's collapse --- but no strategy can bias the forecast or extract an
expectation advantage over the current posterior.
\item[\textbf{(F3b)}] \THMC{Conj.\ \textsf{CHT-POVM} $+$ Pos.\ thermal-time}.
The Colbeck--Renner clause: the terminal PVM is the maximal abelian centraliser
of the tracial state on the cosmological-horizon algebra, so no predictive
extension beyond $\{E_k\}$ exists.
% parent: thm:bell-no-go-evasion, item (4) (l.22992); conj:CHT-POVM (l.21659)
Stated with the correct time index: no variable available strictly before the
terminal limit, and no extension of the framework's statistics, predicts $K$
better than the posterior. (Under terminal record-completeness $K$ \emph{is}
a.s.\ a function of the complete pre-terminal record history --- this note's
first version wrongly denied it; openness is openness-at-every-interior-moment,
not ontological chanciness of the completed block. The compatibilist reading of
this fact is discussed, not asserted, in Section~\ref{sec:notproven}.)
\item[\textbf{(F3c)}] \THM{} under terminal record-completeness (Part-I core;
type-free via the predual theorem); quantitative residual form without it.
Collapse $\pi^\gamma_{s_n}(k)\to\mathbf 1_{\{K=k\}}$ holds only in the terminal
limit,
% parent: thm:collapse (l.1735); ax:complete (l.1654); fd:thm:predual (l.14898)
and for accessible sub-filtrations the residual entropy $H(K\mid\mathcal G)\ge0$
can remain strictly positive: definiteness never precedes the records.
% parent: prop:gen-collapse (l.1893); cor:horizon (l.1954); fd:thm:quantitative (l.14972)
\item[\textbf{(F4a)}] \THM{}; Part-I core, active law. The Born weight is not
invariant under bulk interventions: under the active law the terminal marginal
is the disturbed weight, dependent on the chosen instruments; the change is a
causal influence on a non-local quantity, not a signal, and backward influence
is permitted within the post-selected sub-ensemble.
% parent: remark ``Why p(k) is not invariant'' + rem:gsssbp (l.1493--1540);
% fd:thm:instrument (iv) disturbed Born law (l.14678);
% capacity quantification rem:JLW-retrocausal-capacity (l.1544)
\item[\textbf{(F4b)}] \THM{}; Part-I core only (the Ehrenfest identity is
formulated by the parent in the finite-dimensional setting and no
field-theoretic analogue is claimed here). Outcome-wise conservation of
conditioned expectations and the conditioned Ehrenfest theorem: \emph{between}
interventions, conditioned expectations along the realized branch obey the
ordinary conserved Hamiltonian dynamics; the intervention step itself is the
instrument update.
% parent: thm:conservation (l.1067); prop:ehrenfest (l.2178); instrument update
% fd:thm:instrument (l.14678)
\item[\textbf{(F4c)}] \STRUCT{}$\mid$record layer (RP), per worldline. The
parent's delayed-choice scope statement holds by construction of the
(RP)-admissible protocol class --- a deferred commit/erase decider is outside
the class until its decision becomes a record in the causal past of the readout
--- supported at proposition grade by protocol-counterfactual independence of
the metric outside record cones.
% parent: delayed-choice scope discussion (l.2454--2465); protocol-counterfactual
% independence clause (l.3328--3343)
\emph{Open corner, inherited unshrunk:} the reconciliation of \emph{distinct}
worldlines whose deferred choices are mutually spacelike is the parent's open
multi-worldline corner; the multi-being form of this clause is not claimed.
% parent: closing sentence of the delayed-choice scope passage (l.2463--2465)
\end{itemize}
\end{theorem}

\begin{remark}[Proof discipline]\label{rem:discipline}
Each clause is one of: quoted from a parent theorem at its printed grade (with
label and line anchor in the source comment); imported from the ambient
literature at \ADOPTED{} grade; verified by inspection of the parent axiom set
(\STRUCT{}); or a one-step corollary derived in this note with its proof
sketched in place (\DERIVED{}: the instrument construction in F1c, the optional
sampling step in F3a). No clause is asserted above the grade of its weakest
ingredient; the conjunction is graded in Table~\ref{tab:grades}.
\end{remark}

\begin{remark}[Regime scoping of the conjunction]\label{rem:regimes}
No single regime carries all thirteen clauses at unconditional grade, and the
theorem does not pretend otherwise. In the finite-dimensional Part-I core, the
spine F2a, F2c, F2d, F3a, F3c, F4a, F4b is theorem-grade, but F1a can fail and
F1b/F1c are not available (no Reeh--Schlieder). In the field-theoretic regime,
F1a--F1c hold under the net bundle, F3a/F3c transfer type-free via the parent's
algebra-free predual collapse theorem, two-boundary conditioning itself is
theorem-grade on the Araki branch (the parent's (E4), closed unconditionally),
and the sandwich-state form of F2d rides on \textsf{conj:tomita}, which is
hereby named in the bundle rather than silently consumed. F4b is claimed in the
core only. A single field-theoretic being therefore satisfies the conjunction
at the grades: F1 as printed; F2a structural; F2b conditional; F2c in the core
idealization; F2d on the Araki branch; F3 as printed; F4a--F4b in the core
idealization; F4c per worldline. This is the strongest honest form of the claim; concentration, not
concealment.
\end{remark}

\section{The equality corollary}\label{sec:equality}

\begin{corollary}[Equal possibility of all beings]\label{cor:equal}
At any moment, all beings have the same possibilities of choice, in the
following precise and graded sense.
\begin{itemize}
\item[(i)] \emph{Same options} (grade: F1a's). Their choice sets are identified
by $*$-isomorphism of their diamond algebras: there is no algebraic fact about
what one being \emph{could do} that is not equally a fact about every other.
\item[(ii)] \emph{Same passive law} (grade: \THM{} for worldline-independence;
\ADOPTED{}$\mid$Pos.\ thermal-time bundle for frame-independence). Under the
passive two-boundary law the terminal distribution $p(k)$ is
worldline-independent,
% parent: lem:born (l.754); remark ``Worldline dependence'' (l.1810--1816)
and, for ideal QRF observers, frame-independent --- the frame-independence half
resting on the parent's observer-invariance proposition, itself an adoption
conditional on the thermal-time postulate and the gravity-sector realisation,
and weakening for non-ideal QRFs.
% parent: prop:observer-invariance, tier (A) and its hypothesis line (l.23957)
Under the \emph{active} law the realized terminal marginal is
protocol-dependent --- that is F4a's efficacy, shared symmetrically by all
beings, not an inequality between them.
\item[(iii)] \emph{Same reach} (grade: F1c's). Each being's local operations
reach a dense set of global conditioned states; none is granted a region of
state space inaccessible to another.
\item[(iv)] \emph{Same openness} (grade: F3a unconditional; F3b conditional).
Each being's forecast of the realized branch is bounded by the same martingale
structure, and no being possesses --- or, conditional on the F3b bundle, could
possess --- a better predictor. No being is the universe's oracle.
\item[(v)] \emph{Same time, modulo gauge} (grade: postulate-conditional ---
Pos.\ thermal-time, with the parent's co-scoping residual and
\textsf{conj:tomita} seed behind it; not theorem-grade). Thermal time is
state-dependent only at the inner level; the Connes cocycle is inner, so beings
share a canonical outer flow, their clocks differing exactly by their
thermodynamic state.
% parent: pos:thermal-time (a \begin{postulate}, l.6235); rem:thermal-time-physics,
% item (i) (l.6289)
\end{itemize}
What distinguishes beings is exhaustively indexical: \emph{which} subalgebra
(location), \emph{which} filtration (history), \emph{which} inner cocycle
(thermodynamic state), and --- per (F1b) --- which options the branch's terminal
record forecloses. None of these alters the abstract possibility structure, the
passive law, the reach, or the openness. The weakest-link rule of
Table~\ref{tab:grades} covers every anchor cited by an item, parent postulates
and adopted propositions included, not only F-clauses.
\end{corollary}

\begin{remark}[What equality does \emph{not} say]\label{rem:not-equal}
Equality of possibility is not equality of knowledge or of practical capacity.
Horizon-restricted beings carry strictly positive residual entropy about the
terminal record while others may not;
% parent: cor:horizon (l.1954)
and the magnitude of a being's influence on the realized law depends on its
location and instruments. The corollary asserts equality of the
\emph{possibility structure, passive law, reach, and openness} --- not of
epistemic access, resources, or leverage. Conflating these would overclaim, and
the parent's own horizon-entropy corollary refutes the conflation.
\end{remark}

\begin{remark}[Equality and mutual influence]\label{rem:mutual}
Non-coercion is a statement about marginals and must not be inflated. What is
theorem-protected (F2c) is each being's $k$-summed spacelike marginal
statistics. What is \emph{not} protected --- and is in fact asserted by F4a,
symmetrically for every being --- is the shared terminal law and the per-branch
post-selected statistics: any being's intervention shifts the disturbed weights
$\widetilde p(k)$ that all beings jointly face, and backward influence within a
post-selected sub-ensemble is permitted by the parent.
% parent: rem:gsssbp (l.1519); active-law marginal non-invariance (l.1493--1540)
In a multi-agent world, beings are mutually \emph{influential} at the level of
the branch law while mutually \emph{non-coercive} at the level of marginals ---
the same distinction the parent draws between causal influence and signalling.
Freedom here is freedom among interdependent beings, not isolation.
\end{remark}

\begin{remark}[Why this is freedom-shaped and not merely randomness]\label{rem:not-noise}
Chance alone is not freedom; the conjunction is what carries the content. (F2)
says nothing outside the being forces or biases its options at the level of its
marginal statistics --- not the initial data (measurement independence), not
other beings (no-signalling), not the future boundary (prior-average inertness)
--- while Remark~\ref{rem:mutual} states exactly the sense in which beings
remain mutually influential; (F3) says nothing, inside or outside the
framework, forecasts the resolution better than the being's own posterior; (F1)
says the option set is never degraded by the conditioning except through the
branch's terminal record; (F4) says the choice genuinely reshapes the law of
what is realized, including retro-shaping the post-selected past, lawfully and
causally. A fixed future boundary that is provably non-coercive at the marginal
level, non-predictive, option-preserving, and choice-sensitive is the strongest
freedom-compatible structure a two-boundary theory can offer.
\end{remark}

\section{What is not proven}\label{sec:notproven}

\begin{itemize}
\item \CONCEDED{} \emph{Libertarian free will.} Per realization, the two-boundary
block --- initial data, terminal record $K$, and the record history --- is a
fixed fact, and under terminal record-completeness $K$ is a.s.\ determined by
the complete pre-terminal record history. The framework offers no ontological
``could have done otherwise'' over and above (F2a)'s protocol-externality and
(F3b)'s no-better-predictor clause. The theorem is compatibilist in kind: it
proves the fixity is undetectable in prior average (F2d), non-coercive at the
marginal level (F2), non-predictive from any interior moment (F3), and
option-preserving (F1). It does not dissolve the block; it proves the block
cannot be used against its inhabitants.
\item \CONCEDED{} \emph{A theory of deliberation.} The protocol $\PP_\gamma$ is
an unanalyzed input. The framework does not model how a being chooses; it
proves the choice, however made, enters as an external variable the dynamics
does not weight (F2a) and cannot anticipate (F3, F4c). This is a hospitality
theorem, not a mechanism of agency.
\item \CONCEDED{} \emph{Unconditional status of every clause.} (F2b) and (F3b)
ride on Conj.\ \textsf{CHT-POVM} \emph{and} Pos.\ thermal-time --- two named
points of conditionality, the first being the same one the parent flags for its
no-go evasions; (F1a--F1c) carry the net bundle (III$_1$, Hadamard/RS state,
factoriality, separable predual, split); the field-theoretic sandwich form of
(F2d) rides on \textsf{conj:tomita}; (F4b) is core-only; (F4c) is per-worldline
with the multi-worldline corner open. Failure of \textsf{CHT-POVM} degrades
(F2b) to the parent's symmetry-natural fine-tuning fallback but eliminates
(F3b)'s no-predictive-extension clause outright --- no fallback for it exists in
the parent --- leaving openness supported only by the unconditional martingale
and collapse-timing clauses (F3a, F3c).
% parent: rem:WS-conditional fallback governs the fine-tuning defence only
% (app:wood-spekkens, l.22491--22583)
\item \CONCEDED{} \emph{Derivation of the Born weight.} The equality-of-law
clause inherits the parent's status register: the Born weight is an input, not
a derived quantity.
% parent: rem:born-status (l.787--811); rem:fpf-born (l.1818)
\end{itemize}

\section{Attack ledger}\label{sec:attacks}

Anticipated referee attacks, recorded in the ledger idiom; each with a verdict.
FW-A1 through FW-A7 were drafted with the note; FW-A8 and FW-A9 were added
after the adversarial swarm run, which also forced the re-tensing of FW-A3.

\begin{itemize}
\item[\textbf{FW-A1}] \emph{``Block fatalism: $K$ is fixed, so nothing is
free.''} Verdict: \textsc{contained}, with grades displayed. The fixity is
per-realization and per-protocol: the joint law is a function of the being's
free protocol choice (F4a, unconditional), the boundary's fixity is provably
inert in prior average (F2d, unconditional in the core), and no interior-moment
variable predicts the resolution (F3a unconditional; the no-extension leg F3b
conditional on \textsf{CHT-POVM}$+$thermal-time). The residue --- ``given the
whole block, this block'' --- is a tautology, conceded in
Section~\ref{sec:notproven} and carrying no operational content by (F2d).
\item[\textbf{FW-A2}] \emph{``Equal possibility is vacuous algebra.''} Verdict:
\textsc{refuted}. The isomorphism clause fails in the finite-dimensional core
and holds precisely because local algebras in the field-theoretic regime are
the unique hyperfinite III$_1$ factor. Equality of beings is a contingent,
load-bearing structural fact about the universe's local algebras, not a
notational triviality.
\item[\textbf{FW-A3}] \emph{``Projective terminal effects extinguish options,
and \textsf{CHT-POVM} selects a PVM.''} Verdict: \textsc{conceded-and-scoped},
re-tensed after review. On projective branches faithfulness holds on the
support corner and fails exactly on $\ker L_k$ (F1b). The extinguished options
are those incompatible with the branch's \emph{terminal record} --- a
future-boundary datum, not the being's past. The defence is therefore not the
past-constraint trope (this note's first version used it; the swarm killed it)
but non-coercion: the future-record constraint is block-level conditioning that
is operationally inert in prior average (F2d), undetectable and unusable from
inside, and identical in kind for every being. Constraint by the branch one is
in fact on is what conditioning \emph{means}; what freedom requires --- and
what the theorem delivers --- is that the constraint is not wielded by anything
against anyone. Cross-reference FW-A1.
\item[\textbf{FW-A4}] \emph{``Retrocausal setting-dependence is superdeterminism
by the back door.''} Verdict: \textsc{refuted at grade}. The dependence runs
from settings to past hidden variables within post-selected sub-ensembles ---
not from initial data to settings. Measurement independence at the initial
slice is retained verbatim in the parent; cosmic Bell tests are structurally
blind to future-boundary conditioning. Grade honestly carried: conditional on
\textsf{CHT-POVM}$+$thermal-time, with the parent's lawful-fine-tuning
concession inherited.
\item[\textbf{FW-A5}] \emph{``The `beings' are QRF triples; no agent, no will.''}
Verdict: \textsc{conceded by design}. See the metaphysics fence. The theorem's
content survives the deflation: whatever does the choosing, the framework
does not weight it, cannot coerce it at the marginal level, cannot foresee it,
and does not asymmetrically endow it.
\item[\textbf{FW-A6}] \emph{``Equality is false: a being behind a horizon knows
less.''} Verdict: \textsc{conceded and quarantined}. Epistemic access is
provably unequal (horizon-hidden entropy); the corollary claims equality of
possibility structure, passive law, reach, and openness only
(Remark~\ref{rem:not-equal}).
\item[\textbf{FW-A7}] \emph{``Universal reach is postselection theater: the
success probability can be astronomically small.''} Verdict:
\textsc{conceded-and-scoped}. (F1c) is a possibility statement, graded as such;
it asserts the same dense reach for all beings, not cheap reach for any. The
note's freedom notion nowhere trades on confusing possibility with facility.
\item[\textbf{FW-A8}] \emph{``Beings coerce each other through the shared law:
A's intervention changes the $\widetilde p(k)$ that B faces.''} Verdict:
\textsc{conceded-and-scoped} (Remark~\ref{rem:mutual}). True, symmetric, and
not coercion in the theorem's sense: B's spacelike marginals are
theorem-protected (F2c); the shared branch law is influence-sensitive for
everyone alike (F4a). Mutual influence at the law level with mutual
non-coercion at the marginal level is the framework's precise shape of
interdependent freedom; the note claims nothing stronger.
\item[\textbf{FW-A9}] \emph{``Agents reasoning about agents: Wigner's friend
and Frauchiger--Renner break shared-world reasoning.''} Verdict:
\textsc{delegated at grade}. The parent's Bell-evasion theorem, item (3),
makes nested agents' reasoning chains well-defined relative to the unique
terminal record, blocking the Frauchiger--Renner contradiction ---
\THMC{\textsf{CHT-POVM}$+$thermal-time}, the same bundle as F2b/F3b.
% parent: thm:bell-no-go-evasion, item (3) (l.22992); delegation of the
% measurement-paradox face per prop:observer-invariance tier (C) (l.23957)
A being coherently manipulated as another's measured system is outside the
record-generating protocol class; the equality corollary speaks of
record-generating beings.
\end{itemize}

\section{Grade table}\label{sec:grades}

\begin{table}[h]
\centering\scriptsize\setlength{\tabcolsep}{2.5pt}
\begin{tabular}{@{}lllll@{}}
\toprule
Clause & Content & Grade & Regime & Parent anchor \\
\midrule
F1a & isomorphic choice algebras & \ADOPTED{} $+$ net bundle & QFT & RS-ABL hyp.; ambient \\
F1b & non-extinction of options & \THMC{net} (dense-range; corner o.w.) & QFT & RS-ABL (i)--(iv) \\
F1c & universal reach & \THMC{net} $+$ \DERIVED{} & QFT & RS-ABL (ii) \\
F2a & no weighting of protocol & \STRUCT{} (scoped; adm.\ conditions exist) & both & filtration section \\
F2b & measurement independence & \THMC{\textsf{CHT-POVM}$+$th.-time} & QFT & Bell evasion (1) \\
F2c & marginal non-coercion & \THM{} & core & bulk no-signalling \\
F2d & inert future boundary & \THM{} core; $\mid$\textsf{conj:tomita} QFT & both & prior-average identity \\
F3a & martingale openness & \THM{} $+$ \DERIVED{}; type-free & both & posterior lemma; predual thm. \\
F3b & no better predictor & \THMC{\textsf{CHT-POVM}$+$th.-time} & QFT & Bell evasion (4) \\
F3c & definiteness only with records & \THM{}$\mid$record-completeness; type-free & both & collapse thm.; horizon cor. \\
F4a & efficacy on the branch law & \THM{} (active law) & core & $p(k)$ non-invariance \\
F4b & lawful consequences & \THM{} & core only & conservation; Ehrenfest \\
F4c & causal, not anticipating & \STRUCT{}$\mid$(RP); per worldline & core & delayed-choice scope \\
Cor.~\ref{cor:equal} & equality of beings & weakest-link over \emph{all} cited anchors & --- & incl.\ Pos.\ th.-time, obs.-inv. \\
\bottomrule
\end{tabular}
\caption{Grades carried faithfully; the corollary's weakest-link rule ranges
over every anchor its items cite --- F-clauses, parent postulates, and adopted
propositions alike. Unconditional spine (each in its stated regime): F2a
(scoped), F2c, F2d (core), F3a, F3c ($\mid$record-completeness), F4a, F4b
(core). F1b/F1c are theorem-grade under the net bundle; nothing in this note is
unconditional in a regime it is not stated for.}
\label{tab:grades}
\end{table}

\section*{Ambient references (import list)}
\noindent Cited from standard knowledge; primary-source verification pending
per house rule before any import into the parent.
\begin{itemize}
\item U. Haagerup, \emph{Connes' bicentralizer problem and uniqueness of the
injective factor of type III$_1$}, Acta Math.\ 158 (1987) 95--148. [F1a]
\item A. Connes, \emph{Classification of injective factors}, Ann.\ Math.\ 104
(1976) 73--115. [F1a]
\item D. Buchholz, C. D'Antoni, K. Fredenhagen, \emph{The universal structure
of local algebras}, Commun.\ Math.\ Phys.\ 111 (1987) 123--135. [F1a: split,
and the hyperfinite-III$_1$-factor-up-to-center reduction that makes the
factoriality hypothesis explicit]
\item C. List, \emph{Free will, determinism, and the possibility of doing
otherwise}, No\^us 48 (2014) 156--178. [compatibilist framing, fence]
\end{itemize}

\end{document}
